\section{Additional Exercises}

\subsection{Where Connection Coefficients Appear}

\begin{exercise}\label{ex:parallel-1}
Determine the coefficients of the Riemann tensor with respect to a chart $(U,x)$ in terms of the connection coefficient functions.
\end{exercise}
\hyperref[sol:parallel-1]{\small \textit{Solution on p.~\pageref*{sol:parallel-1}}}

\begin{exercise}\label{ex:parallel-2}
Does a one-dimensional manifold with connection have curvature? Why?
\end{exercise}
\hyperref[sol:parallel-2]{\small \textit{Solution on p.~\pageref*{sol:parallel-2}}}

\subsection{The Round Sphere}

\textit{There is a question about finding the components of Riem for given $\Gamma$s. I am not going to write it here to save myself time, but I recommend at least watching the video (available \href{https://gravity-and-light.herokuapp.com/tutorials}{here}, for some reason this video is not on YouTube) for a worked calculation.}

\subsection{How Not To Define Parallel Transport}

\begin{exercise}\label{ex:parallel-3}
H\"{a}nschen defines two vectors $X\in T_p\cM$ and $Y\in T_q\cM$ as parallel if
\bse
    X^i_{(x)} = Y^i_{(x)}
\ese
with respect to some chart $(U,x)$ whose domain $U$ contains both $p$ and $q$.

Prove that this notion of parallelity is ill-defined!
\end{exercise}
\hyperref[sol:parallel-3]{\small \textit{Solution on p.~\pageref*{sol:parallel-3}}}

\textit{There is then a diagram to show how the above definition of parallelity fails for parallel (as defined in the lecture) vectors around a circle when considered in the Cartesian chart and the polar chart.

To save myself time I have not drawn the diagrams here, but if you can't picture the images in your head, again go see the video!}
